\centerline{\bf \Huge Малая теорема Ферма}
\centerline{\bf \Huge }

\begin{flushright}
        \scriptsize{\href{https://www.youtube.com/watch?v=xwvKcWa2vLM&t=10s}{А в каком кольце?}}
\end{flushright}

\problem{1}
Число $p>5$ - простое. Докажите, что число, составленное из $p-1$ единицы, делится на $p$.

\problem{2}
Докажите, что для любого целого числа $a$

a) $a^5-a$ делится на 30 ;

b) $a^{17}-a$ делится на 510 ;

c) $a^{11}-a$ делится на 66 .

\problem{3}
Найдите остаток от деления

(a) $2^{100}$ на 101 ;

(b) $7^{102}$ на 101 ;

(c) $8^{900}$ на 29 .

\problem{4}
Пусть $a$ взаимно просто с 1001. Докажите, что $a^{3000}-1$ делится на 1001.

\problem{5}
Существует ли степень тройки, оканчивающаяся на 0001 ?

\problem{6}
Докажите, что $2^{3^k}+1$ делится на $3^{k+1}$.

\problem{7}
Пусть $m$ и $n$ натуральные числа.

a) Докажите, что $\varphi(n) \cdot \varphi(m)=\varphi($ НОД $(m, n)) \cdot \varphi($ НОК $(m, n))$.

б)$^{*}$ Докажите \textbf{усиление теоремы Эйлера}. 

Положим $L(m)=\operatorname{HOK}\left(\varphi\left(p_1^{t_1}\right), \ldots, \varphi\left(p_k^{t_k}\right)\right)$. Тогда, если $(a, m)=1$, то $a^{L(m)}-1$ делится на $m$.

\problem{8$^{**}$}
Докажите великую теорему Ферма в случае когда $n$ достаточно большое простое число (больше чем остальные переменные из уравнения).